\chapter{GL(1): Dirichlet $L$-functions}
\label{chap:gl1}

The GL(1) case of Deligne's conjecture is the classical arithmeticity of critical
values of Dirichlet $L$-functions: the Artin motive of a primitive Dirichlet character
$\chi$ modulo $N$. This chapter gives a complete proof.

Throughout, $N \geq 1$, $\chi$ is a primitive Dirichlet character modulo $N$ with
values in $\C$, and $\delta \in \{0,1\}$ is defined by $\chi(-1) = (-1)^{\delta}$.
Mathlib records the dichotomy as \lean{DirichletCharacter.Even} ($\delta = 0$) and
\lean{DirichletCharacter.Odd} ($\delta = 1$), exhaustive by
\lean{DirichletCharacter.even_or_odd}.

\section{The data}

\begin{definition}[Field of values]
  \label{def:valueField}
  \lean{Deligne.GL1.valueField}
  \leanok
  $\Q(\chi) \subseteq \C$ is the intermediate field $\Q$ adjoin the values of $\chi$. In Lean,
  \lean{IntermediateField.adjoin ℚ (Set.range fun a => χ a)}.
\end{definition}

\begin{lemma}[The field of values is algebraic]
  \label{lem:valueField-algebraic}
  \lean{Deligne.GL1.isAlgebraic\_of\_mem\_valueField}
  \leanok
  \uses{def:valueField, lem:values-algebraic}
  Every element of $\Q(\chi)$ is algebraic over $\Q$, which is requirement (ix) of
  Definition~\ref{def:pkg}.
\end{lemma}
\begin{proof}
  \leanok
  Every value $\chi(a)$ is an algebraic number (Lemma~\ref{lem:values-algebraic}), and a field
  generated over $\Q$ by algebraic elements is algebraic
  (\lean{IntermediateField.isAlgebraic\_adjoin\_iff\_isAlgebraic}).
\end{proof}

\begin{definition}[Galois conjugate]
  \label{def:galoisConj}
  \lean{Deligne.GL1.galoisConj}
  \leanok
  For $\sigma \in \operatorname{Gal}(\C/\Q)$, set $\chi^{\sigma} = \sigma \circ \chi$,
  which is again a Dirichlet character modulo $N$ (\lean{MulChar.ringHomComp}).
\end{definition}

\begin{definition}[Critical integers]
  \label{def:gl1-critical}
  \lean{Deligne.GL1.IsCritical}
  \leanok
  An integer $n$ is \emph{critical} for $\chi$ if either
  \begin{itemize}
    \item $n \geq 1$ and $n \equiv \delta \pmod 2$, or
    \item $n \leq 0$ and $n \not\equiv \delta \pmod 2$.
  \end{itemize}
\end{definition}

\begin{definition}[Period]
  \label{def:gl1-period}
  \lean{Deligne.GL1.period}
  \leanok
  \[
    c^{+}(M(\chi)(n)) \;=\;
    \begin{cases}
      \tau(\chi)\,(2\pi i)^{n}, & n \geq 1,\\
      1, & n \leq 0,
    \end{cases}
  \]
  where $\tau(\chi) = \operatorname{gaussSum} \chi\ \lean{ZMod.stdAddChar}$.
\end{definition}

The branch is not an artefact of bookkeeping: it is exactly where $c^{+}$ and $c^{-}$
exchange roles as $n$ crosses the centre of the critical strip, and it is the
structure that becomes the $\omega^{\pm}$ bookkeeping for GL(2).

\begin{remark}[Sanity check against $\zeta$]
  \label{rem:zeta-sanity}
  \uses{def:gl1-critical, def:gl1-period}
  For $\chi = \mathbf{1}$ (so $\delta = 0$ and $L(\mathbf{1},s) = \zeta(s)$),
  Definition~\ref{def:gl1-critical} gives critical integers $2,4,6,\dots$ and
  $-1,-3,-5,\dots$. This is correct: $\zeta(2k)$ is a rational multiple of
  $\pi^{2k}$ and $\zeta(1-2k) = -B_{2k}/2k$ is rational, while $\zeta(3)$ and the
  trivial zeros $\zeta(-2k)$ are correctly excluded.

  The period normalisation is pinned down by three numerical checks, each of which
  the final theorem must reproduce:
  \begin{itemize}
    \item $\chi$ the odd character mod $4$, $n=1$: $L(\chi,1) = \pi/4$, $\tau(\chi) = 2i$,
      so $c^{+} = 2i \cdot 2\pi i = -4\pi$ and the ratio is $-1/16 \in \Q$;
    \item $\chi = \mathbf{1}$, $n=2$: $\zeta(2) = \pi^2/6$, $c^{+} = (2\pi i)^2 = -4\pi^2$,
      ratio $-1/24 \in \Q$;
    \item $\chi = \mathbf{1}$, $n=4$: $\zeta(4) = \pi^4/90$, $c^{+} = 16\pi^4$,
      ratio $1/1440 \in \Q$.
  \end{itemize}
  Note $\tau(\chi) \notin \Q(\chi)$ in general --- it lives in $\Q(\chi,\zeta_N)$ ---
  so the period carries genuine transcendental content and is not a normalisation
  that could be absorbed into the field of values.
\end{remark}

\begin{lemma}[Critical set of the Riemann zeta function]
  \label{lem:gl1-crit-spec}
  \lean{Deligne.GL1.isCritical_one_iff}
  \leanok
  \uses{def:gl1-critical}
  This is the obligation of Remark~\ref{rem:crit-spec}, and it must be a check against
  something \emph{independently} specified --- an unfolding of
  Definition~\ref{def:gl1-critical} against itself would be a tautology and would
  certify nothing. We therefore specialise to $\chi = \mathbf{1}$, where the answer is
  known in advance, and assert
  \[ \mathrm{IsCritical}(\mathbf{1}, n) \iff (2 \leq n \wedge 2 \mid n) \vee (n \leq -1 \wedge 2 \nmid n). \]
  The right-hand side mentions neither $\mathrm{IsCritical}$ nor
  \lean{Deligne.Pkg}, and is the critical set of $\zeta$ as classically recorded.
\end{lemma}
\begin{proof}
  \leanok
  The trivial character is even, so $\delta = 0$ and the two branches of
  Definition~\ref{def:gl1-critical} read $n \geq 1 \wedge 2 \mid n$ and
  $n \leq 0 \wedge 2 \nmid n$. Since no odd $n$ is $0$ and no even $n \geq 1$ is $1$,
  these are the two displayed conditions.
\end{proof}

\begin{lemma}[Numerical normalisation checks]
  \label{lem:gl1-numeric}
  \lean{Deligne.GL1.normalizedValue_modOne_two, Deligne.GL1.normalizedValue_modOne_four}
  \uses{def:gl1-period, def:normalizedValue, lem:gaussSum-modOne}
  The second half of the obligation of Remark~\ref{rem:crit-spec}: the period of
  Definition~\ref{def:gl1-period} is checked against known values, so that a period
  wrong by a nonzero constant is detected. Assert, as Lean statements,
  \[ \widetilde{L}(\mathbf{1},2) = -\tfrac{1}{24}, \qquad
     \widetilde{L}(\mathbf{1},4) = \tfrac{1}{1440}, \qquad
     \widetilde{L}(\chi_{-4},1) = -\tfrac{1}{16}, \]
  where $\chi_{-4}$ is the odd character modulo $4$.
\end{lemma}
\begin{proof}
  The first two follow from \lean{riemannZeta\_two} and \lean{riemannZeta\_four}
  together with $\tau(\mathbf{1}) = 1$ (Lemma~\ref{lem:gaussSum-modOne}), which
  collapses the period to $c^{+} = (2\pi i)^{n}$. The third needs
  $L(\chi_{-4},1) = \pi/4$ and $\tau(\chi_{-4}) = 2i$; the Gauss sum is a two-term
  finite computation, and the $L$-value is the Leibniz series, so this one may need a
  separate sub-blueprint if mathlib does not supply it directly. Recorded here as a
  statement; see Remark~\ref{rem:numeric-status}.
\end{proof}

\begin{remark}[Status of the numerical checks]
  \label{rem:numeric-status}
  The two $\zeta$ checks are formalized, as
  \lean{Deligne.GL1.normalizedValue\_modOne\_two} and
  \lean{Deligne.GL1.normalizedValue\_modOne\_four}; the $\chi_{-4}$ check is not yet, and has no
  Lean statement in the project. The $\chi_{-4}$ check is the one
  that pins the Gauss sum factor, since $\tau(\mathbf{1}) = 1$ makes the $\zeta$
  checks blind to it, and it is therefore the one worth the effort. If
  $L(\chi_{-4},1) = \pi/4$ is not directly available it follows from
  Theorem~\ref{thm:gl1-pos} itself --- but then it is no longer an independent check,
  and should be replaced by one that is.
\end{remark}

The next three lemmas are the obligations Definition~\ref{def:pkg} and \S\ref{sec:guards}
impose. They sit at different depths, and the difference is the subject of
Remark~\ref{rem:not-a-pkg}: Lemma~\ref{lem:gl1-period-ne-zero} discharges requirement (x)
and Lemma~\ref{lem:values-algebraic} requirement (ix), both fields of the structure, so without
them Definition~\ref{def:gl1-pkg} does not typecheck at all; and
Lemma~\ref{lem:gl1-nondegenerate} discharges
Definition~\ref{def:hasCriticalInteger}, a predicate proved about the package afterwards, so
without it the package still exists but says nothing at any object.

\begin{lemma}[Nondegeneracy: a critical integer exists]
  \label{lem:gl1-nondegenerate}
  \lean{Deligne.GL1.exists\_isCritical, Deligne.GL1.hasCriticalInteger\_pkg}
  \leanok
  \uses{def:gl1-critical, def:hasCriticalInteger}
  Every $\chi$ has a critical integer; equivalently, every object of
  Definition~\ref{def:gl1-pkg} is nondegenerate in the sense of
  Definition~\ref{def:hasCriticalInteger}.
\end{lemma}
\begin{proof}
  \leanok
  Take $n = 1$ if $\chi$ is odd and $n = 2$ if $\chi$ is even; both are positive and
  congruent to $\delta$ modulo $2$. Exhaustiveness is
  \lean{DirichletCharacter.even_or_odd}.
\end{proof}

\begin{lemma}[Requirement (x): the period does not vanish]
  \label{lem:gl1-period-ne-zero}
  \lean{Deligne.GL1.period\_ne\_zero}
  \leanok
  \uses{def:gl1-period, def:pkg, cor:gaussSum-ne-zero}
  Let $\chi$ be primitive modulo $N \geq 1$. Then $c^{+}(M(\chi)(n)) \neq 0$ for
  \emph{every} $n \in \Z$, critical or not.
\end{lemma}
\begin{proof}
  \leanok
  Split on the branch of Definition~\ref{def:gl1-period}. For $n \leq 0$ the period is
  $1$. For $n \geq 1$ it is $\tau(\chi)\,(2\pi i)^{n}$, a product of two nonzero
  factors: $\tau(\chi) \neq 0$ is Corollary~\ref{cor:gaussSum-ne-zero}, which is where
  primitivity is used and is the only substantive input; and
  $(2\pi i)^{n} \neq 0$ since $\pi \neq 0$ and $i \neq 0$, an integer power of a
  nonzero complex number being nonzero.

  The statement is deliberately stronger than requirement (x) demands --- it drops the
  criticality hypothesis, which the proof never uses. Making it unconditional keeps the
  lemma independent of Definition~\ref{def:gl1-critical}, so a later change to the
  critical set cannot break the package's definition.

  Note the shape of the dependency: primitivity is a hypothesis on \emph{objects of the
  package}, and the package is indexed by primitive characters precisely so that this proof
  has the hypothesis available at every object. That is the reason for the subtype in
  Definition~\ref{def:gl1-pkg}, and it is forced by requirement (x), not chosen for
  convenience.
\end{proof}

\begin{lemma}[The values of $\chi$ are algebraic]
  \label{lem:values-algebraic}
  \lean{DirichletCharacter.isAlgebraic\_apply}
  \leanok
  Every value $\chi(a)$ of a Dirichlet character modulo $N \geq 1$ is algebraic over $\Q$.
\end{lemma}
\begin{proof}
  \leanok
  If $a$ is not a unit then $\chi(a)=0$. If it is, then
  $\chi^{\,\#(\Z/N)^{\times}} = 1$, so $\chi(a)$ is a root of unity, hence a root of
  $X^{\#(\Z/N)^{\times}} - 1$.
\end{proof}

It is far from sharp --- $\Q(\chi)$ is a subfield of $\Q(\zeta_{N})$, so a number field --- but
it is all that is needed; see Remark~\ref{rem:valueField-bounded}.

\section{The non-positive half}
\label{sec:gl1-negative}

This is the engine of the whole case. At a critical $n \leq 0$ the period is $1$, so
the assertion is that $L(\chi,n)$ is \emph{itself} an element of $\Q(\chi)$. The
positive half is then deduced from this by the functional equation, because
$n \mapsto 1-n$ exchanges the two halves of the critical set. So this half cannot be
postponed: it is what the other half is proved from.

\begin{definition}[Generalised Bernoulli number]
  \label{def:genBernoulli}
  \lean{DirichletCharacter.genBernoulli}
  \leanok
  For $\chi$ modulo $N$ valued in a field $L$ of characteristic zero,
  \[ B_{k,\chi} \;=\; N^{k-1} \sum_{a=1}^{N} \chi(a)\, B_k\!\left(\tfrac{a}{N}\right), \]
  where $B_k$ is the $k$-th Bernoulli polynomial (\lean{Polynomial.bernoulli}).
  Washington, \emph{Introduction to Cyclotomic Fields}, \S4.1, Prop.~4.1. The range
  $1 \leq a \leq N$ matters: it makes the trivial-character case reduce to
  $B_k(1) = \lean{bernoulli'}\ k$.
\end{definition}

\begin{lemma}[Generalised Bernoulli numbers are rational combinations of character values]
  \label{lem:genBernoulli-ratCast}
  \lean{DirichletCharacter.genBernoulli\_eq\_ratCast\_mul\_sum}
  \leanok
  \uses{def:genBernoulli}
  Let $L$ be a field of characteristic zero, $\iota : \Q \to L$ its canonical embedding, and
  $\chi$ a Dirichlet character modulo $N$ valued in $L$. Then
  \[ B_{k,\chi} \;=\; \iota(r_{N,k}) \sum_{a=0}^{N-1} \chi(a+1)\, \iota(b_{k,a}),
     \qquad r_{N,k} = N^{k-1} \in \Q, \qquad b_{k,a} = B_k\!\left(\tfrac{a+1}{N}\right) \in \Q. \]
  That is: every factor of Definition~\ref{def:genBernoulli} other than the values of $\chi$
  themselves is the image of an explicit rational number.
\end{lemma}
\begin{proof}
  \leanok
  Definition~\ref{def:genBernoulli} already has this shape, so the only point needing an
  argument is that the Bernoulli polynomial may be evaluated in $\Q$ \emph{before} being mapped
  to $L$ rather than after. That is
  \lean{Polynomial.aeval\_algebraMap\_apply\_eq\_algebraMap\_eval} applied to the
  $\Q$-algebra $L$: evaluating $\iota_{*}B_k$ at $\iota(q)$ is $\operatorname{aeval}$, which
  commutes with the structure map, giving
  $\operatorname{eval}_{\iota(q)}(\iota_{*}B_k) = \iota(\operatorname{eval}_q B_k)$. The
  argument $(a+1)/N$ is itself $\iota$ of a rational, as is $N^{k-1}$.

  This single fact is what both Lemma~\ref{lem:genBernoulli-mem} and
  Lemma~\ref{lem:genBernoulli-galois} rest on, which is why it is separated out. Once the
  non-character factors are exhibited as rationals, the first is closed by the fact that every
  subfield of $\C$ contains $\Q$, and the second by the fact that every automorphism of $\C$
  fixes $\Q$ pointwise. Neither needs to know anything further about Bernoulli polynomials.
\end{proof}

\begin{theorem}[Values at non-positive integers]
  \label{thm:LFunction-neg-nat}
  \lean{PadicLFunctions.LFunction\_neg\_nat}
  \leanok
  \uses{def:genBernoulli, lem:hurwitz-neg-nat-Ioo}
  For every $k \in \N$,
  \[ L(\chi,-k) \;=\; -\,\frac{B_{k+1,\chi}}{k+1}. \]
\end{theorem}
\begin{proof}
  \uses{lem:hurwitz-neg-nat-Ioo}
  Split on $N$.

  \emph{Level one.} $\chi = \mathbf{1}$ by \lean{DirichletCharacter.level\_one}, and
  \lean{ZMod.LFunction\_modOne\_eq} identifies $L(\mathbf{1},s)$ with $\zeta(s)$.
  Conclude with \lean{riemannZeta\_neg\_nat\_eq\_bernoulli'} and the evaluation
  $B_k(1) = \lean{bernoulli'}\ k$ of Definition~\ref{def:genBernoulli}.

  \emph{Level $N > 1$.} Expand using
  \[ \lean{ZMod.LFunction}\ \Phi\ s \;=\; N^{-s} \sum_{j \in \Z/N} \Phi(j)\, \zeta_H\!\left(\tfrac{j}{N}, s\right), \]
  which is the definition of \lean{ZMod.LFunction}. Since $\chi$ is a character of
  $(\Z/N)^{\times}$ and $0$ is not a unit for $N > 1$, we have $\chi(0) = 0$
  (\lean{MulChar.map\_nonunit}), so the $j = 0$ term drops and the sum runs over
  $j/N \in (0,1)$ strictly. Apply Lemma~\ref{lem:hurwitz-neg-nat-Ioo} at $s = -k$ to
  each remaining term and collect: the result is exactly
  $-B_{k+1,\chi}/(k+1)$ after rewriting $N^{k}$ against the $N^{k+1-1}$ of
  Definition~\ref{def:genBernoulli}.
\end{proof}

\begin{lemma}[Hurwitz zeta at non-positive integers, interior form]
  \label{lem:hurwitz-neg-nat-Ioo}
  \lean{PadicLFunctions.hurwitzZeta\_neg\_nat\_of\_mem\_Ioo}
  \leanok
  For every $k \in \N$ and every real $x$ with $0 < x < 1$,
  \[ \zeta_H(x,-k) \;=\; \frac{-1}{k+1}\, B_{k+1}(x). \]
\end{lemma}
\begin{proof}
  \uses{lem:hurwitzZetaOdd-zero}
  For $k \neq 0$ this is mathlib's \lean{HurwitzZeta.hurwitzZeta\_neg\_nat}, whose
  hypotheses are $k \neq 0$ and $x \in [0,1]$ --- both satisfied.

  For $k = 0$ mathlib's lemma does not apply, and its docstring records this as an
  open TODO. Decompose $\zeta_H = \zeta_H^{\mathrm{even}} + \zeta_H^{\mathrm{odd}}$.
  By \lean{HurwitzZeta.hurwitzZetaEven\_apply\_zero},
  \[ \zeta_H^{\mathrm{even}}(x,0) = \begin{cases} -1/2 & x \equiv 0, \\ 0 & \text{otherwise},\end{cases} \]
  and $x \in (0,1)$ gives $x \not\equiv 0$ in $\R/\Z$, so the even part vanishes.
  Lemma~\ref{lem:hurwitzZetaOdd-zero} supplies the odd part, and
  $B_1(x) = x - 1/2$ (\lean{Polynomial.bernoulli\_one}) finishes.

  The restriction to the \emph{open} interval is essential and not a technicality:
  at $x = 0$ the identity genuinely fails, since
  $\zeta_H(0,0) = \zeta(0) = -1/2 \neq -B_1(0) = 1/2$. This is the jump of the
  sawtooth, and it is why Theorem~\ref{thm:LFunction-neg-nat} needs $\chi(0) = 0$.
\end{proof}

\begin{lemma}[The sawtooth boundary value]
  \label{lem:hurwitzZetaOdd-zero}
  \lean{PadicLFunctions.hurwitzZetaOdd\_apply\_zero\_of\_mem\_Ioo}
  \leanok
  For real $x$ with $0 < x < 1$, $\zeta_H^{\mathrm{odd}}(x,0) = \tfrac12 - x$.
\end{lemma}
\begin{proof}
  Apply \lean{HurwitzZeta.hurwitzZetaOdd\_one\_sub} at $s = 1$ (its hypothesis
  $\forall n \in \N,\ s \neq -n$ holds), giving
  \[ \zeta_H^{\mathrm{odd}}(x,0) \;=\; 2 (2\pi)^{-1} \Gamma(1) \sin(\pi/2)\, Z_{\sin}(x,1)
     \;=\; \pi^{-1} Z_{\sin}(x,1). \]
  So the claim reduces to the boundary value of the sine zeta function,
  $Z_{\sin}(x,1) = \pi(\tfrac12 - x)$, which is the classical sawtooth Fourier series
  $\sum_{n \geq 1} \sin(2\pi n x)/n = \pi(\tfrac12 - x)$ on $(0,1)$. This is not in
  mathlib: the Dirichlet series defining $Z_{\sin}$ converges only conditionally at
  $s = 1$, so it must be reached by an Abel limit
  (\lean{Mathlib.Analysis.Complex.AbelLimit}) together with a uniform bound on the
  partial sums of $\sum \sin(2\pi n x)$.

  \textbf{Source.} That argument is not carried out here. This lemma, and
  Lemma~\ref{lem:hurwitz-neg-nat-Ioo} and Theorem~\ref{thm:LFunction-neg-nat} above it,
  are taken \emph{sorry-free} from AINTLIB
  (\url{https://github.com/CBirkbeck/AINTLIB}), as
  \lean{sinZeta\_one\_eq\_boundary},
  \lean{hurwitzZetaOdd\_apply\_zero\_of\_mem\_Ioo} and
  \lean{hurwitzZeta\_neg\_nat\_of\_mem\_Ioo} in
  \path{projects/PadicLFunctions/PadicLFunctions/Interpolation/Sawtooth.lean}
  (617 lines), and as \lean{LFunction\_neg\_nat} in the adjacent
  \path{GenBernoulliComplex.lean} (74 lines). \emph{Sawtooth.lean} imports only
  mathlib plus AINTLIB's \path{Common/Analysis/DirichletBounds.lean}, so the needed
  slice is four files. All of it has been built in this workspace, which is what the
  \leanok\ above records; see Remark~\ref{rem:aintlib}.
\end{proof}

\begin{remark}[On depending on AINTLIB]
  \label{rem:aintlib}
  \emph{Resolved.} This project and AINTLIB now agree on
  \lean{leanprover/lean4:v4.34.0-rc2} and on mathlib \texttt{e4b72ca0d01c}, so the two
  dependencies resolve to a single mathlib rather than conflicting.

  The open risk this remark used to record was that AINTLIB's own \path{lakefile.toml}
  marks the \lean{PadicLFunctions} library as ``need a mathlib bump to compile \dots\
  built on demand'' and excludes it from \lean{defaultTargets}, so it was not
  known-green even against AINTLIB's own pin. It has since been built in this
  workspace: \path{PadicLFunctions.Interpolation.GenBernoulliComplex} and its
  transitive closure, including \path{Sawtooth.lean}, compile without error. The
  \leanok\ marks on Theorem~\ref{thm:LFunction-neg-nat},
  Lemma~\ref{lem:hurwitz-neg-nat-Ioo} and Lemma~\ref{lem:hurwitzZetaOdd-zero} record
  that, and rest on an actual build here rather than on AINTLIB's say-so.

  The fallback is therefore no longer needed, but it is worth keeping on record: were
  the dependency ever dropped, only Lemma~\ref{lem:hurwitzZetaOdd-zero} would have to
  be reproved in-project from the Abel-limit argument above, and the gap would be
  confined to $n = 0$ for odd $\chi$.
\end{remark}

\begin{lemma}[Rationality of generalised Bernoulli numbers]
  \label{lem:genBernoulli-mem}
  \lean{Deligne.GL1.genBernoulli\_mem\_valueField}
  \leanok
  \uses{def:genBernoulli, def:valueField, lem:genBernoulli-ratCast}
  $B_{k,\chi} \in \Q(\chi)$.
\end{lemma}
\begin{proof}
  \leanok
  Immediate from Definition~\ref{def:genBernoulli}: $N^{k-1}$ is rational, the
  Bernoulli polynomial has rational coefficients and is evaluated at the rational
  points $a/N$, each $\chi(a)$ lies in $\Q(\chi)$ by construction, and
  $\Q(\chi)$ is a subfield of $\C$ hence contains $\Q$ and is closed under the finite
  sum and products involved.
\end{proof}

\begin{lemma}[Galois equivariance of generalised Bernoulli numbers]
  \label{lem:genBernoulli-galois}
  \lean{Deligne.GL1.genBernoulli\_galoisConj}
  \leanok
  \uses{def:genBernoulli, def:galoisConj, lem:genBernoulli-ratCast}
  $\sigma(B_{k,\chi}) = B_{k,\chi^{\sigma}}$ for every $\sigma \in \operatorname{Gal}(\C/\Q)$.
\end{lemma}
\begin{proof}
  \leanok
  $\sigma$ is a ring homomorphism, so it commutes with the finite sum and the
  products in Definition~\ref{def:genBernoulli}. It fixes $\Q$ pointwise, hence fixes
  $N^{k-1}$ and each value $B_k(a/N)$ of a rational polynomial at a rational point.
  The only remaining factors are the $\chi(a)$, and $\sigma(\chi(a)) = \chi^{\sigma}(a)$
  by Definition~\ref{def:galoisConj}.
\end{proof}

\begin{theorem}[Non-positive half]
  \label{thm:gl1-nonpos}
  \lean{Deligne.GL1.isArithmetic\_of\_nonpos, Deligne.GL1.isEquivariant\_of\_nonpos,
    Deligne.GL1.normalizedValue\_of\_nonpos, Deligne.GL1.normalizedValue\_neg\_natCast}
  \leanok
  \uses{thm:LFunction-neg-nat, lem:genBernoulli-mem, lem:genBernoulli-galois, def:gl1-period}
  Let $n \leq 0$ be critical for $\chi$. Then $L(\chi,n)/c^{+}(M(\chi)(n)) \in \Q(\chi)$,
  and for every $\sigma \in \operatorname{Gal}(\C/\Q)$,
  $\sigma(L(\chi,n)) = L(\chi^{\sigma},n)$.
\end{theorem}
\begin{proof}
  \leanok
  \uses{thm:LFunction-neg-nat, lem:genBernoulli-mem, lem:genBernoulli-galois}
  Write $n = -k$ with $k \in \N$. By Definition~\ref{def:gl1-period} the period is
  $1$, so the normalised value is $L(\chi,-k)$ itself. Algebraicity is
  Theorem~\ref{thm:LFunction-neg-nat} followed by
  Lemma~\ref{lem:genBernoulli-mem} (and closure of $\Q(\chi)$ under division by the
  nonzero rational $k+1$). Equivariance is Theorem~\ref{thm:LFunction-neg-nat} on both
  sides with Lemma~\ref{lem:genBernoulli-galois} in between; the period is $1$ on both
  sides and $\sigma(1) = 1$.

  Note that criticality is not used here: Theorem~\ref{thm:LFunction-neg-nat} holds at
  every non-positive integer, so the conclusion does too. (At the non-critical
  non-positive integers the value is in fact $0$, by the parity vanishing
  $B_{k,\chi} = 0$ for $k \not\equiv \delta \pmod 2$ --- but that vanishing is not
  stated in this blueprint and nothing below depends on it, so it is not claimed
  here.)
\end{proof}

\section{The positive half}
\label{sec:gl1-positive}

This is the half where the period is a genuine transcendental, and so the half that
actually exercises the framework. It is a corollary of \S\ref{sec:gl1-negative} and
mathlib's functional equation, not an independent argument.

\begin{lemma}[When two characters share a critical integer]
  \label{lem:gl1-common-critical}
  \lean{Deligne.GL1.exists\_isCritical\_and\_isCritical\_iff}
  \leanok
  \uses{def:gl1-critical}
  Two Dirichlet characters $\chi_1, \chi_2$ (of any moduli) have a common critical integer if
  and only if they have the same parity.
\end{lemma}
\begin{proof}
  \leanok
  By Definition~\ref{def:gl1-critical} the critical set of an even character is
  $\{n \geq 2 \text{ even}\} \cup \{n \leq -1 \text{ odd}\}$ and that of an odd character is
  $\{n \geq 1 \text{ odd}\} \cup \{n \leq 0 \text{ even}\}$. Same parity: take $n = 2$ if both
  are even and $n = 1$ if both are odd. Opposite parity: the two sets are disjoint, since on
  $n \geq 1$ they demand opposite parities of $n$ and on $n \leq 0$ they likewise do.

  This is the obstruction of Remark~\ref{rem:sum-obstruction}: it is why
  Definition~\ref{def:hasCriticalInteger} is not inherited by Definition~\ref{def:sum}, and
  why $\zeta_K$ for $K$ imaginary quadratic has no critical integers.
\end{proof}

\begin{lemma}[Criticality transfers across the functional equation]
  \label{lem:crit-one-sub}
  \lean{Deligne.GL1.isCritical\_one\_sub\_iff}
  \leanok
  \uses{def:gl1-critical, lem:even-inv}
  Let $n \geq 1$. Then $n$ is critical for $\chi$ if and only if $1-n \leq 0$ is
  critical for $\chi^{-1}$.
\end{lemma}
\begin{proof}
  \leanok
  $\chi^{-1}$ has the same parity as $\chi$ (Lemma~\ref{lem:even-inv}). Writing $n = 2r + \delta$ with $r \geq 0$, we get
  $1 - n = 1 - 2r - \delta$, so $1-n \equiv 1 - \delta \not\equiv \delta \pmod 2$,
  which is exactly the non-positive branch of Definition~\ref{def:gl1-critical}.
\end{proof}

\begin{lemma}[Value field is inversion-invariant]
  \label{lem:valueField-inv}
  \lean{Deligne.GL1.valueField\_inv}
  \leanok
  \uses{def:valueField}
  $\Q(\chi^{-1}) = \Q(\chi)$.
\end{lemma}
\begin{proof}
  \leanok
  On units $\chi^{-1}(a) = \chi(a)^{-1}$, and off the units both characters vanish. So
  every value of $\chi^{-1}$ is either $0$ or the inverse of a value of $\chi$, and
  $\Q(\chi)$ --- being a subfield of $\C$ --- contains both. Hence
  $\operatorname{range}\chi^{-1} \subseteq \Q(\chi)$ and, by minimality of
  \lean{IntermediateField.adjoin}, $\Q(\chi^{-1}) \subseteq \Q(\chi)$. The reverse inclusion is
  the same argument with $\chi$ and $\chi^{-1}$ exchanged, using $(\chi^{-1})^{-1} = \chi$.

  (No appeal to complex conjugation is needed. It is true that the values are roots of
  unity and hence that $\chi^{-1}(a) = \overline{\chi(a)}$, but that requires
  $|\chi(a)| = 1$ and buys nothing here.)
\end{proof}

\begin{lemma}[The Gamma factor does not vanish at the mirror point]
  \label{lem:gamma-ne-zero}
  \lean{Deligne.GL1.gammaFactor\_neg\_natCast\_ne\_zero}
  \leanok
  \uses{def:gl1-critical, lem:gammaR-odd-ne-zero}
  Let $k \in \N$ and suppose $k+1$ is critical for $\chi$. Then $\gamma(\chi, -k) \neq 0$.
\end{lemma}
\begin{proof}
  \leanok
  Write $\delta = 0$ if $\chi$ is even and $\delta = 1$ if $\chi$ is odd, so that
  $\gamma(\chi,s) = \Gamma_{\R}(s+\delta)$. Criticality at $n = k+1 \geq 1$ says
  $n \equiv \delta \pmod 2$, i.e.\ $k+1+\delta$ is even, so $k+\delta$ is odd and therefore
  so is $-k+\delta$. Now apply Lemma~\ref{lem:gammaR-odd-ne-zero}.

  Note that the companion statement $\gamma(\chi,k+1) \neq 0$
  (Lemma~\ref{lem:gammaFactor-pos-ne-zero}) needs no criticality hypothesis at all, which is
  why the two live in different chapters: only this one is about the critical set.
\end{proof}

\begin{lemma}[Root number times Gamma quotient]
  \label{lem:gamma-quotient}
  \lean{Deligne.GL1.rootNumber\_mul\_gammaFactor\_quotient}
  \leanok
  \uses{def:gl1-critical, lem:gammaR-refl-even, lem:gammaR-refl-odd, lem:gammaC-natCast,
    lem:cos-natCast-mul-pi}
  Let $k \in \N$ and suppose $k+1$ is critical for $\chi$. Then
  \[ W(\chi)\,\frac{\gamma(\chi,-k)}{\gamma(\chi,k+1)}
     \;=\; \frac{\tau(\chi)}{N^{1/2}} \cdot
           \frac{\chi(-1)\,(2\pi i)^{k+1}}{2 \cdot k!}, \]
  where $W(\chi) = \lean{rootNumber}\ \chi = \tau(\chi)\, i^{-\delta} N^{-1/2}$ and
  $\delta \in \{0,1\}$ is the parity of $\chi$.
\end{lemma}
\begin{proof}
  \leanok
  This is the archimedean heart of the positive half, and it is deliberately stated as one
  identity rather than two. The factor $i^{-\delta}$ inside $W(\chi)$ and the sign produced by
  the reflection formula \emph{each} depend on the parity of $\chi$; their product does not.
  Separating them would force mathlib's conditional
  $i^{(\text{if } \chi \text{ even then } 0 \text{ else } 1)}$ into the statement, and would
  then require it to be cancelled again at the point of use.

  Write $n = k+1$ and split on parity. In both cases the hypothesis of
  Lemmas~\ref{lem:gammaR-refl-even} and~\ref{lem:gammaR-refl-odd} at $s = n$ holds, since
  $n \geq 1$ is not of the form $-j$ with $j \in \N$.

  \emph{$\chi$ even ($\delta = 0$).} Criticality makes $n$ even, say $k = 2r+1$ and
  $n = 2r+2$. Here $\gamma(\chi,s) = \Gamma_{\R}(s)$, and $-k = 1-n$, so
  Lemma~\ref{lem:gammaR-refl-even} gives
  $\gamma(\chi,-k)/\gamma(\chi,n) = (\Gamma_{\C}(n)\cos(\pi n/2))^{-1}$. Now
  $\cos(\pi n/2) = \cos((r+1)\pi) = (-1)^{r+1}$ by Lemma~\ref{lem:cos-natCast-mul-pi}, and
  $\Gamma_{\C}(n)^{-1} = 2^{k}\pi^{n}/k!$ by Lemma~\ref{lem:gammaC-natCast}. On the right,
  $\chi(-1) = 1$ and $i^{n} = i^{2r+2} = (-1)^{r+1}$; both sides are
  $\tau(\chi) N^{-1/2} (-1)^{r+1} 2^{k}\pi^{n}/k!$.

  \emph{$\chi$ odd ($\delta = 1$).} Criticality makes $n$ odd, say $k = 2r$ and $n = 2r+1$.
  Here $\gamma(\chi,s) = \Gamma_{\R}(s+1)$, and $-k+1 = 2-n$, so
  Lemma~\ref{lem:gammaR-refl-odd} gives
  $\gamma(\chi,-k)/\gamma(\chi,n) = (\Gamma_{\C}(n)\sin(\pi n/2))^{-1}$, with
  $\sin(\pi n/2) = \cos(r\pi) = (-1)^{r}$, again by
  Lemma~\ref{lem:cos-natCast-mul-pi} after $\sin(x + \pi/2) = \cos x$. On the right, $\chi(-1) = -1$ and
  $i^{n} = i^{2r+1} = (-1)^{r} i$, while on the left $W(\chi)$ carries $i^{-1} = -i$; both
  sides are $-i\,\tau(\chi) N^{-1/2} (-1)^{r} 2^{k}\pi^{n}/k!$.
\end{proof}

\begin{proposition}[The functional equation at a positive critical integer]
  \label{prop:pos-fe}
  \lean{Deligne.GL1.LFunction\_natCast\_add\_one}
  \leanok
  \uses{lem:gamma-quotient, lem:gamma-ne-zero, lem:gammaFactor-inv,
    lem:gammaFactor-pos-ne-zero}
  Let $\chi$ be primitive modulo $N$, let $k \in \N$ and suppose $k+1$ is critical for
  $\chi$. Then
  \[ L(\chi, k+1) \;=\; \tau(\chi)\,(2\pi i)^{k+1} \cdot
       \frac{\chi(-1)}{2\,N^{k+1}\,k!} \cdot L(\chi^{-1}, -k). \]
\end{proposition}
\begin{proof}
  \leanok
  \uses{lem:gamma-quotient, lem:gamma-ne-zero}
  Mathlib's functional equation
  \lean{DirichletCharacter.IsPrimitive.completedLFunction\_one\_sub} reads, for primitive
  $\chi$,
  \[ \Lambda(\chi, 1-s) \;=\; N^{s-1/2}\, W(\chi)\, \Lambda(\chi^{-1}, s). \]
  Put $s = -k$, so $1-s = k+1$. Convert both completed functions to ordinary ones with
  \lean{DirichletCharacter.LFunction\_eq\_completed\_div\_gammaFactor}, read in the direction
  $\Lambda = L\cdot\gamma$: this is licensed by
  Lemma~\ref{lem:gammaFactor-pos-ne-zero} at $s = k+1$ and by Lemma~\ref{lem:gamma-ne-zero}
  at $s = -k$, the latter transported from $\chi^{-1}$ to $\chi$ by
  Lemma~\ref{lem:gammaFactor-inv}. That lemma's side condition ($s \neq 0$ or $N \neq 1$)
  holds at $s = k+1$ because $k+1 \neq 0$, and at $s = -k$ because it can fail only when
  $k = 0$ and $N = 1$; but the character modulo $1$ is even, so criticality of $n = 1$ would
  demand that $1$ be even. Rearranging,
  \[ L(\chi,k+1) \;=\; N^{-k-1/2}\, W(\chi)\,
       \frac{\gamma(\chi^{-1},-k)}{\gamma(\chi,k+1)}\, L(\chi^{-1},-k), \]
  and Lemma~\ref{lem:gammaFactor-inv} followed by Lemma~\ref{lem:gamma-quotient} evaluates the
  middle factor. The two half-powers of $N$ combine,
  $N^{-k-1/2}\cdot N^{-1/2} = N^{-(k+1)}$, an integral power, which is the first of the two
  cancellations of Remark~\ref{rem:pos-cancellations}.
\end{proof}

\begin{proposition}[Reduction of the positive half]
  \label{prop:pos-reduction}
  \lean{Deligne.GL1.normalizedValue\_natCast\_add\_one}
  \leanok
  \uses{def:gl1-period, prop:pos-fe, lem:gl1-period-ne-zero}
  Let $\chi$ be primitive modulo $N$, let $n = k+1 \geq 1$ be critical for $\chi$. Then
  \[ \frac{L(\chi,n)}{c^{+}(M(\chi)(n))} \;=\;
     \frac{\chi(-1)}{2\,N^{n}\,(n-1)!}\; L(\chi^{-1},1-n). \]
  In particular the right-hand side is $\chi(-1)$ times a \emph{rational} multiple of
  $L(\chi^{-1},1-n)$, and the Gauss sum has cancelled.
\end{proposition}
\begin{proof}
  \leanok
  \uses{prop:pos-fe}
  By Definition~\ref{def:gl1-period} the period at $n \geq 1$ is
  $c^{+} = \tau(\chi)(2\pi i)^{n}$, which is exactly the first factor of
  Proposition~\ref{prop:pos-fe} and is nonzero by
  Lemma~\ref{lem:gl1-period-ne-zero}. Divide.
\end{proof}

\begin{remark}[The scalar, and two cancellations]
  \label{rem:pos-cancellations}
  The closed form of Proposition~\ref{prop:pos-reduction} is simpler than the shape in which
  such a formula is usually quoted, and the simplification is not cosmetic --- it is what
  makes both halves of Theorem~\ref{thm:gl1-pos} short. Running the computation through the
  half-integer values of $\Gamma$ instead produces
  \[ \frac{(-1)^{m} q_{n,\delta}}{(2N)^{n}}, \qquad
     q_{n,\delta} = \frac{(-4)^{r} r!}{(2r)!\,\Gamma(r+\delta)}, \qquad n = 2r+\delta, \;
     m = r+\delta, \]
  and one checks, in each parity separately, that this equals
  $\chi(-1)/(2N^{n}(n-1)!)$: for $\delta = 0$ it is
  $(-1)^{2r} 2^{2r-1}/((2r-1)!\,2^{2r}N^{n})$, and for $\delta = 1$ it is
  $-2^{2r}/((2r)!\,2^{2r+1}N^{n})$. The sign is $(-1)^{\delta} = \chi(-1)$, and using
  $\chi(-1)$ rather than $(-1)^{\delta}$ is what removes parity bookkeeping from the proofs
  of Theorem~\ref{thm:gl1-pos} and Theorem~\ref{thm:gl1-pos-equiv}: $\chi(-1)$ is visibly a
  value of $\chi$, hence in $\Q(\chi)$, and $\sigma(\chi(-1)) = \chi^{\sigma}(-1)$ is the
  definition of the Galois action.

  Two cancellations make the answer land in $\Q(\chi)$ rather than a larger field. The
  $\sqrt{N}$ inside $W(\chi)$ cancels against $N^{s-1/2}$ at $s = 1-n$, leaving an integral
  power of $N$. The $i^{-\delta}$ inside $W(\chi)$ cancels against the $i^{n}$ coming from
  $(2\pi i)^{n}$, and this works \emph{only} because $n \equiv \delta \pmod 2$. Criticality is
  doing real work here; it is also the reason the period of Definition~\ref{def:gl1-period}
  needs no separate power of $i$.

  Finally, the formula can be checked against the two values already computed independently
  in Lemma~\ref{lem:gl1-numeric}. At $N = 1$ and $n = 2$ it gives
  $\zeta(-1)/2 = -1/24$, and at $n = 4$ it gives $\zeta(-3)/12 = 1/1440$; both agree, using
  $\zeta(-1) = -1/12$ and $\zeta(-3) = 1/120$. These are the only two independent numerical
  anchors the development currently has, and the closed form was verified against both before
  it was formalised.
\end{remark}

\begin{theorem}[Positive half: algebraicity]
  \label{thm:gl1-pos}
  \lean{Deligne.GL1.isArithmetic\_of\_pos}
  \leanok
  \uses{prop:pos-reduction, thm:gl1-nonpos, lem:valueField-inv, lem:primitive-inv}
  Let $n \geq 1$ be critical for $\chi$. Then $L(\chi,n)/c^{+}(M(\chi)(n)) \in \Q(\chi)$.
\end{theorem}
\begin{proof}
  \leanok
  \uses{prop:pos-reduction, thm:gl1-nonpos, lem:valueField-inv, lem:primitive-inv}
  Write $n = k+1$. By Lemma~\ref{lem:primitive-inv}, $\chi^{-1}$ is primitive, so
  Theorem~\ref{thm:gl1-nonpos} applies to it at $-k$ and gives
  $L(\chi^{-1},-k) \in \Q(\chi^{-1})$, which is $\Q(\chi)$ by
  Lemma~\ref{lem:valueField-inv}. Proposition~\ref{prop:pos-reduction} exhibits the
  normalised value as $\chi(-1)$ times a rational multiple of it; $\chi(-1)$ is a value of
  $\chi$ and so lies in $\Q(\chi)$ by definition, and $\Q(\chi)$ contains $\Q$.

  One step of the expected argument is absent, and its absence is worth recording.
  Theorem~\ref{thm:gl1-nonpos} is proved for \emph{every} non-positive integer, not only the
  critical ones --- the generalised Bernoulli number $B_{k+1,\chi}$ lies in $\Q(\chi)$
  whatever the parity of $k$. So the criticality of $1-n$ for $\chi^{-1}$ is never needed, and
  Lemma~\ref{lem:crit-one-sub} is not invoked here. It is kept as a statement about the
  critical set in its own right, not as a step of this proof.
\end{proof}

\begin{theorem}[Positive half: equivariance]
  \label{thm:gl1-pos-equiv}
  \lean{Deligne.GL1.isEquivariant\_of\_pos}
  \leanok
  \uses{prop:pos-reduction, thm:gl1-nonpos, lem:primitive-inv, lem:pkg-galoisAction-coe,
    def:galoisConj}
  Let $n \geq 1$ be critical for $\chi$ and let $\sigma \in \operatorname{Gal}(\C/\Q)$. Then
  $\sigma\bigl(L(\chi,n)/c^{+}(M(\chi)(n))\bigr)$ is the corresponding normalised value of
  $\chi^{\sigma}$.
\end{theorem}
\begin{proof}
  \leanok
  \uses{prop:pos-reduction, thm:gl1-nonpos, lem:pkg-galoisAction-coe, lem:gl1-galois-transport}
  Write $n = k+1$. Criticality of $n$ for $\chi^{\sigma}$ is
  Lemma~\ref{lem:gl1-galois-transport}, so
  Proposition~\ref{prop:pos-reduction} applies to $\chi^{\sigma}$ as well, with the
  \emph{same} scalar: it depends only on $N$, $n$ and $\chi(-1)$, and
  $\sigma(\chi(-1)) = \chi^{\sigma}(-1)$. Apply $\sigma$ to
  Proposition~\ref{prop:pos-reduction}. The rational factor $1/(2N^{n}(n-1)!)$ is fixed by
  $\sigma$ because $\sigma$ fixes $\Q$ pointwise, so what remains is
  $\sigma(L(\chi^{-1},-k)) = L((\chi^{\sigma})^{-1},-k)$. That is the equivariance half of
  Theorem~\ref{thm:gl1-nonpos} applied to $\chi^{-1}$, together with
  $(\chi^{-1})^{\sigma} = (\chi^{\sigma})^{-1}$ --- both are $a \mapsto \sigma(\chi(a))^{-1}$,
  and in Lean this is mathlib's \lean{MulChar.ringHomComp\_inv}. As in
  Theorem~\ref{thm:gl1-pos}, no criticality of $-k$ is needed.
\end{proof}

\section{Assembly}

\begin{definition}[Galois conjugation as a group action]
  \label{def:galoisAction}
  \lean{Deligne.GL1.galoisAction}
  \leanok
  \uses{def:galoisConj, lem:primitive-galois}
  For $\sigma \in \operatorname{Gal}(\C/\Q)$, the map $\chi \mapsto \chi^{\sigma}$ is a
  permutation of the primitive Dirichlet characters modulo $N$, with inverse
  $\chi \mapsto \chi^{\sigma^{-1}}$, and $\sigma \mapsto (\chi \mapsto \chi^{\sigma})$ is a
  group homomorphism $\operatorname{Gal}(\C/\Q) \to \operatorname{Sym}(\{\chi\})$.
\end{definition}
\begin{proof}
  \leanok
  Well defined by Lemma~\ref{lem:primitive-galois}, and the two round trips are
  $\sigma^{-1}(\sigma(z)) = z$ and $\sigma(\sigma^{-1}(z)) = z$ applied to each value;
  $(\chi^{\sigma})^{\tau} = \chi^{\sigma\tau}$ and $\chi^{1} = \chi$ valuewise. This is the
  form clause (vi) of Definition~\ref{def:pkg} requires, since the Galois action there is a
  homomorphism into $\operatorname{Sym}(\Omega)$ rather than a bare map.
\end{proof}

\begin{lemma}[Primitivity is Galois-stable]
  \label{lem:primitive-galois}
  \lean{Deligne.GL1.isPrimitive\_galoisConj}
  \leanok
  \uses{def:galoisConj, lem:primitive-ringHomComp}
  If $\chi$ is primitive then so is $\chi^{\sigma}$.
\end{lemma}
\begin{proof}
  \leanok
  $\chi^{\sigma} = \sigma \circ \chi$ and a ring automorphism is injective, so this is
  Lemma~\ref{lem:primitive-ringHomComp}. That lemma is the substance: mathlib has no
  result on the conductor under post-composition, and without one this step --- on which
  the well-definedness of the whole package depends --- has no proof.
\end{proof}

\begin{lemma}[Galois conjugation preserves parity and transports the field of values]
  \label{lem:gl1-galois-transport}
  \lean{Deligne.GL1.even\_galoisConj, Deligne.GL1.isCritical\_galoisConj\_iff,
    Deligne.GL1.valueField\_galoisConj}
  \leanok
  \uses{def:galoisConj, def:gl1-critical, def:valueField}
  For every $\sigma \in \operatorname{Gal}(\C/\Q)$, $\chi^{\sigma}$ has the parity of $\chi$,
  hence the same critical set, and $\Q(\chi^{\sigma}) = \sigma(\Q(\chi))$. These are
  requirements (vii) and (viii) of Definition~\ref{def:pkg} for the package below.
\end{lemma}
\begin{proof}
  \leanok
  $\sigma$ fixes $1$ and is injective, so $\sigma(\chi(-1)) = 1$ exactly when
  $\chi(-1) = 1$; criticality depends on $\chi$ only through its parity. For the field,
  $\sigma$ maps the generating set $\operatorname{range}\chi$ of $\Q(\chi)$ onto the
  generating set of $\Q(\chi^{\sigma})$, so $\Q(\chi^{\sigma}) = \sigma(\Q(\chi))$ by
  \lean{IntermediateField.adjoin\_map}.
\end{proof}

\begin{definition}[The GL(1) package]
  \label{def:gl1-pkg}
  \lean{Deligne.GL1.pkg}
  \leanok
  \uses{def:pkg, def:valueField, def:gl1-critical, def:gl1-period, def:galoisAction,
    lem:primitive-galois, lem:gl1-galois-transport, lem:valueField-algebraic,
    lem:gl1-nondegenerate, lem:gl1-period-ne-zero}
  The Deligne package on the type of primitive Dirichlet characters modulo $N$, with
  $L$ the analytic $L$-function \lean{DirichletCharacter.LFunction}, value field
  $\Q(\chi)$ (Definition~\ref{def:valueField}, \lean{Deligne.GL1.valueField\_pkg}), period
  Definition~\ref{def:gl1-period}, and Galois action Definition~\ref{def:galoisConj}
  (well defined on primitive characters by Lemma~\ref{lem:primitive-galois}). Its weight is $0$
  --- the Artin motive of a Dirichlet character is pure of weight $0$, so the functional equation
  exchanges $s$ and $1-s$ --- and its Gamma shifts are $\{0\}$ for $\chi$ even and $\{1\}$ for
  $\chi$ odd, a single $\Gamma_{\R}$ because the Hodge type is $(0,0)$. The archimedean factor
  that results is mathlib's own \lean{DirichletCharacter.gammaFactor}
  (Lemma~\ref{lem:gl1-gammaFactor-pkg}), and the critical set is computed from it --- see
  Lemma~\ref{lem:gl1-crit-pkg}. Of the requirements of Definition~\ref{def:pkg}, (vii) and
  (viii) are Lemma~\ref{lem:gl1-galois-transport}, (ix) is Lemma~\ref{lem:valueField-algebraic}
  and (x) is Lemma~\ref{lem:gl1-period-ne-zero}; nondegeneracy is
  Lemma~\ref{lem:gl1-nondegenerate}, proved about the package rather than built into it.
\end{definition}

\begin{lemma}[The package's Gamma factor is mathlib's]
  \label{lem:gl1-gammaFactor-pkg}
  \lean{Deligne.GL1.gammaFactor\_pkg, Deligne.GL1.weight\_pkg}
  \leanok
  \uses{def:gl1-pkg}
  For every primitive $\chi$ modulo $N$ and every $s \in \C$, the archimedean factor
  $\prod_{a} \Gamma_{\R}(s+a)$ determined by the shifts of Definition~\ref{def:gl1-pkg} is
  \lean{DirichletCharacter.gammaFactor} $\chi\, s$, and the package's weight is $0$.
\end{lemma}
\begin{proof}
  \leanok
  Split on the parity of $\chi$: the multiset of shifts is a singleton, so the product is
  $\Gamma_{\R}(s+0) = \Gamma_{\R}(s)$ for $\chi$ even and $\Gamma_{\R}(s+1)$ for $\chi$ odd,
  which is mathlib's definition on the nose. This is what makes the check of
  Remark~\ref{rem:crit-spec} a comparison against an external object rather than against
  something this development chose.
\end{proof}

\begin{lemma}[The critical set is the one computed from the Gamma factor]
  \label{lem:gl1-crit-pkg}
  \lean{Deligne.GL1.isCritical\_pkg\_iff}
  \leanok
  \uses{def:gl1-pkg, def:gl1-critical, def:isCritical, lem:gl1-gammaFactor-pkg}
  For every primitive $\chi$ modulo $N$ and every $n \in \Z$, the framework's criticality
  predicate of Definition~\ref{def:isCritical}, computed from the Gamma shifts and reflected
  about $w+1-n = 1-n$, holds at $n$ if and only if Definition~\ref{def:gl1-critical} does.
\end{lemma}
\begin{proof}
  \leanok
  Split on the parity of $\chi$, so that $\gamma(\chi,s) = \Gamma_{\R}(s+\delta)$ with $\delta$
  the parity. By \lean{Complex.Gammaℝ\_eq\_zero\_iff}, $\Gamma_{\R}(m) = 0$ for $m \in \Z$ exactly
  when $m \leq 0$ and $m$ is even. For $\chi$ even the two conditions are that $n$ is not a
  non-positive even integer and $1-n$ is not one either --- that is, $n \notin \{0,-2,-4,\dots\}$
  and $n \notin \{1,3,5,\dots\}$, leaving the even integers $\geq 2$ and the odd integers
  $\leq -1$, which is Definition~\ref{def:gl1-critical} for an even character. For $\chi$ odd the
  same computation at $s+1$ leaves the odd integers $\geq 1$ and the even integers $\leq 0$. Each
  branch is then integer arithmetic.

  This is the obligation of Remark~\ref{rem:crit-spec}, and it is worth saying what has changed:
  Definition~\ref{def:gl1-critical} used to \emph{be} the package's critical set by fiat, so a
  characterisation theorem could only restate it. Now the package's critical set comes from
  mathlib's Gamma factor and this lemma is a genuine comparison, which either holds or does not.
\end{proof}

Coherence needs no separate theorem: criticality is Galois-stable for every package
(Lemma~\ref{lem:isCritical-galois}), and what coherence used to require of GL(1) is now
Lemma~\ref{lem:gl1-galois-transport}, consumed by Definition~\ref{def:gl1-pkg} itself.

\begin{lemma}[The Galois action on the package, unfolded]
  \label{lem:pkg-galoisAction-coe}
  \lean{Deligne.GL1.coe\_galoisAction\_pkg}
  \leanok
  \uses{def:gl1-pkg, def:galoisConj}
  The underlying character of $\sigma \cdot M$ is $\chi^{\sigma}$, where $\chi$ is the
  underlying character of $M$.
\end{lemma}
\begin{proof}
  \leanok
  True by definition. It is recorded (and marked \lean{simp}) because the index type of the
  package is the \emph{subtype} of primitive characters: the Galois action produces an element
  of that subtype packaged with its primitivity proof, and its first component is therefore
  only definitionally, not syntactically, equal to $\chi^{\sigma}$. Every rewrite that pushes
  $\sigma$ through the package needs this bridge.
\end{proof}

\begin{lemma}[Algebraicity, both halves]
  \label{lem:gl1-arithmetic-pkg}
  \lean{Deligne.GL1.isArithmetic\_pkg}
  \leanok
  \uses{thm:gl1-pos, thm:gl1-nonpos, def:gl1-pkg}
  Every primitive $\chi$ modulo $N$ satisfies algebraicity for the package of
  Definition~\ref{def:gl1-pkg}.
\end{lemma}
\begin{proof}
  \leanok
  \uses{thm:gl1-pos, thm:gl1-nonpos}
  Case split on the sign of $n$: Theorem~\ref{thm:gl1-nonpos} for $n \leq 0$ and
  Theorem~\ref{thm:gl1-pos} for $n \geq 1$.
\end{proof}

\begin{lemma}[Equivariance, both halves]
  \label{lem:gl1-equivariant-pkg}
  \lean{Deligne.GL1.isEquivariant\_pkg}
  \leanok
  \uses{thm:gl1-pos-equiv, thm:gl1-nonpos, def:gl1-pkg}
  Every primitive $\chi$ modulo $N$ satisfies equivariance for the package of
  Definition~\ref{def:gl1-pkg}.
\end{lemma}
\begin{proof}
  \leanok
  \uses{thm:gl1-pos-equiv, thm:gl1-nonpos}
  The same case split, with Theorem~\ref{thm:gl1-pos-equiv} on the positive side.
\end{proof}

\begin{theorem}[The critical values are nonzero]
  \label{thm:gl1-nonvanishing}
  \lean{Deligne.GL1.normalizedValue\_ne\_zero}
  \leanok
  \uses{def:gl1-pkg, def:gl1-critical, lem:gl1-period-ne-zero, lem:gl1-nondegenerate}
  For every primitive $\chi$ modulo $N$ and every critical $n \geq 1$,
  $L(\chi,n)/c^{+}(M(\chi)(n)) \neq 0$.
\end{theorem}
\begin{proof}
  \leanok
  The period is nonzero by Lemma~\ref{lem:gl1-period-ne-zero}, so it suffices that
  $L(\chi,n) \neq 0$. This is mathlib's
  \lean{DirichletCharacter.LFunction\_ne\_zero\_of\_one\_le\_re}, which excludes only the
  trivial character at $s = 1$; and that case is never critical, since the character modulo
  $1$ is even and Definition~\ref{def:gl1-critical} then forces $n$ to be even. For $N > 1$ a
  primitive character is nontrivial, because the trivial character has conductor $1$.

  This is what Remark~\ref{rem:numeric-status} was reaching for and did not achieve. The
  guards of \S\ref{sec:guards} rule out the two degenerate \emph{shapes} --- no critical
  integers, and a vanishing period --- but neither prevents an instance whose normalised
  values all happen to be $0$, which would satisfy both halves of the conjecture while
  asserting nothing. This theorem closes that for GL(1), uniformly in the conductor. Note it
  covers only $n \geq 1$; at the non-positive critical integers the value is a generalised
  Bernoulli number and its non-vanishing is a separate question, not needed here because
  Lemma~\ref{lem:gl1-nondegenerate} supplies a positive critical integer for every $\chi$.
\end{proof}

\begin{theorem}[Deligne's conjecture for GL(1)]
  \label{thm:gl1-main}
  \lean{Deligne.GL1.conjecture}
  \leanok
  \uses{lem:gl1-arithmetic-pkg, lem:gl1-equivariant-pkg, thm:gl1-pos,
    thm:gl1-pos-equiv, thm:gl1-nonpos, def:gl1-pkg,
    lem:gl1-nondegenerate, lem:gl1-crit-spec}
  The GL(1) package satisfies Deligne's conjecture: for every primitive Dirichlet
  character $\chi$ modulo $N$ and every critical integer $n$, the normalised value
  $L(\chi,n)/c^{+}(M(\chi)(n))$ lies in $\Q(\chi)$, and $\operatorname{Gal}(\C/\Q)$
  carries it to the corresponding value of $\chi^{\sigma}$. Moreover its critical set is the
  explicit one of Definition~\ref{def:gl1-critical}. Nondegeneracy is
  Lemma~\ref{lem:gl1-nondegenerate}: unlike the non-vanishing of the period, which is a field of
  Definition~\ref{def:pkg} and so is carried by the package itself, it is a separate claim
  about the package --- see Remark~\ref{rem:not-a-pkg}.
\end{theorem}
\begin{proof}
  \leanok
  \uses{lem:gl1-arithmetic-pkg, lem:gl1-equivariant-pkg}
  The conjunction of Lemma~\ref{lem:gl1-arithmetic-pkg} and
  Lemma~\ref{lem:gl1-equivariant-pkg}. The supplementary claim is
  Lemma~\ref{lem:gl1-crit-spec}.
\end{proof}
